\fig{fv:e6:rolling}{向后摩擦时平动功率为负、转动功率为正，无滑动时两项恰好相消}{\begin{tikzpicture}

\draw[world](-1.8,0)--(3,0);\draw[world](0,1.3)circle(1.3);\fill(0,0)circle(.05)node[below]{$P$};
\draw[flow](0,1.3)--(1.8,1.3)node[right]{$V_C>0$};\draw[->](.65,2)arc(50:-75:.95);\node at(1.7,.4){$\omega<0$};
\draw[force](0,0)--(-1.4,0)node[above]{$f<0$};\node[align=left]at(5,1.4){$P_{\text{平}}=fV_C<0$\\$P_{\text{转}}=fR\omega>0$\\$V_C+R\omega=0$};
\end{tikzpicture}}
图\ref{fv:e6:rolling}保持逆时针角速度为正，因此向右纯滚动时$\omega<0$。摩擦向左时$f<0$，平动功率$fV_C$为负，绕质心转动功率$fR\omega$为正，二者相加等于实际作用点的功率$f(V_C+R\omega)=0$。若地面本身运动，应改用相对地面的滑动速度判断摩擦及耗散，不能照搬固定地面的零功结论。
